\documentclass[11pt]{article}

\usepackage[utf8]{inputenc}
\usepackage{geometry}
\geometry{margin=1.0in}
\usepackage{xcolor}
\usepackage{hyperref}
\usepackage{booktabs}
\usepackage{amsmath}
\usepackage{amssymb}
\usepackage{tfrupee}
\usepackage{longtable}

\hypersetup{
    colorlinks=true,
    linkcolor=blue,
    citecolor=blue,
    urlcolor=blue
}

\definecolor{commentcolor}{rgb}{0.1, 0.5, 0.1}
\definecolor{updatecolor}{rgb}{0.0, 0.0, 0.8}

\newcommand{\comment}[1]{\vspace{12pt}\noindent\colorbox{gray!10}{\parbox{\dimexpr\textwidth-2\fboxsep\relax}{\color{commentcolor}\itshape #1}}\vspace{6pt}}
\newcommand{\response}{\vspace{6pt}\noindent\textbf{Response: }}
\newcommand{\update}[1]{\textcolor{updatecolor}{#1}}

\title{\textbf{Response to Reviewer Comments}}
\date{}

\begin{document}

\maketitle

\noindent\textbf{Manuscript ID}: SBR-D-26-00443 \\
\textbf{Title}: AI-Driven E-commerce in the Industry 5.0 Era: A Parallel-Empirical Framework, Decision Support, and Methodological Caution

\section*{Response to the Editor's and Reviewers' Comments}
We express our deepest gratitude to the Editor, Reviewer \#1, and Reviewer \#2 for their thorough evaluation and constructive feedback on our manuscript. We are especially encouraged by Reviewer \#2's recommendation for publication and Reviewer \#1's detailed audit of internal consistency and statistical reporting. 

Below, we provide point-by-point responses detailing all the revisions made.

\subsection*{Point-by-Point Responses to Reviewer \#1}

\comment{\textbf{Comment 1:} \\
The abstract is too lengthy and overloaded with numerical results. It currently reads as a compressed statistical report rather than a focused summary of the research problem, approach, main insight, and contribution. The abstract should not reproduce the exact review count, transaction-row count, holdout accuracy, negative recall, Brier scores, ASIN percentage, two Gini coefficients, Kruskal-Wallis statistic, p value, and effect size. These details belong in the methodology and results. It is sufficient to state that the transformer model performed best, lightweight models remained competitive, the datasets could not be linked, sales were concentrated, and marketplace MRP differences were practically negligible. The abstract should be reduced to approximately 150-200 words and should avoid presenting a sequence of numbers. It also contains the typographical error "utlining."}

\response
We agree with the reviewer and thank them for this constructive suggestion. We have thoroughly rewritten and streamlined the Abstract, condensing it down to about 174 words. We removed all dense numerical listings (exact review counts, holdout accuracy percentages, Brier scores, Gini numbers, $H$-statistics, $p$-values, and effect sizes) and eliminated the typographical error ("utlining"). The revised Abstract now provides a clear, high-level narrative focusing on the core research problem, theoretical framework, parallel findings, and DPDP governance implications:

\begin{quote}
\update{``India's e-commerce ecosystem is evolving rapidly under Artificial Intelligence (AI), digital transformation, and the human-centric pillars of Industry 5.0, contributing to the Viksit Bharat@2047 agenda. Grounded in the Resource-Based View (RBV), Technology-Organization-Environment (TOE), and Technology Acceptance Model (TAM), this study presents a parallel-empirical framework and methodological caution for analysing consumer sentiment and transactional sales performance. Evaluating seven sentiment classifiers on a large product review corpus shows that transformer-based deep learning achieves superior predictive accuracy, while lightweight classical models provide competitive, sustainable, and explainable alternatives suitable for resource-constrained enterprises. Crucially, a data quality linkage audit demonstrates that public review corpora and transactional sales records are disjoint, highlighting a key methodological warning against assuming seamless secondary data integration in digital retail research. Consequently, we demonstrate independent parallel analytics streams—sentiment model benchmarking and operational sales distribution analytics—which reveal significant Gini revenue inequality alongside competitive product tail distributions and negligible cross-marketplace price dispersion. Finally, we discuss how these empirical insights translate into a conceptual enterprise decision-support framework and data governance guidelines under the Digital Personal Data Protection (DPDP) Act of 2023.''}
\end{quote}

\comment{\textbf{Comment 2:} \\
The hypotheses remain weakly connected to the theories. H1a concerns comparative classifier accuracy and calibration. H1b concerns class-balancing procedures. H2a concerns concentration indices. H2b concerns MRP distributions. None follows directly from RBV, TOE, TAM, Stakeholder Theory, or Institutional Theory. These are technical expectations rather than theoretically derived hypotheses.}

\response
We thank the reviewer for this insightful comment regarding theoretical framing. We acknowledge that organizational theories (RBV, TOE, TAM) do not directly predict specific machine learning performance metrics or algorithm-level outputs (such as DistilBERT accuracy or SMOTE recall rates). Instead, these theories establish why AI analytics capabilities and transparent data governance are strategically and organizationally vital for firm competitiveness and regulatory compliance.

To make this distinction clear, we have updated Section 1 (Introduction) to explicitly state that the overarching theoretical frameworks (RBV, TOE, TAM) motivate the organizational necessity for AI analytics assets, data quality audits, and user trust, whereas Hypotheses H1a–H2b serve as empirical implementation tests evaluating competing technical models and operational market distributions. The revised introductory narrative in Section 1 reads as follows:

\begin{quote}
\update{``The theoretical frameworks (RBV, TOE, TAM) establish why AI analytics capabilities and transparent data governance are strategically vital for firm competitiveness and regulatory compliance, whereas our hypotheses (H1a–H2b) empirically evaluate competing technical models and operational market distributions.''}
\end{quote}

\comment{\textbf{Comment 3:} \\
The cross-validation comparison is also not equivalent across models. DistilBERT cross-validation was conducted on a subsample of only 15,000 reviews and trained for one epoch, while its holdout model used more than 446,000 training cases and three epochs. Classical models appear to use the larger corpus. These results do not represent comparable ten-fold validation under a common protocol. The lower DistilBERT cross-validation performance may be an artefact of reduced training data and undertraining rather than instability of the architecture.}

\response
We agree with the reviewer's technical assessment. Conducting 10-fold cross-validation for DistilBERT on a subsample of 15,000 reviews trained for a single epoch was an intentional computational concession to keep cross-validation tractable on local hardware and adhere to Green AI energy considerations. In contrast, classical models (SVM and TF-IDF+LR) were cross-validated on the full dataset partition.

We have updated Sections 3.3, 4.2, and 7.1 (Limitations) to state explicitly that the DistilBERT cross-validation score should not be interpreted as a directly comparable 10-fold benchmark under an identical protocol. We clarify that DistilBERT's cross-validation performance ($93.77\% \pm 0.50\%$) reflects the sensitivity of transformer architectures to training data volume and convergence epochs, whereas its holdout model ($97.69\%$), trained on 446,941 samples over three epochs, demonstrates its true capacity when fully converged.

\comment{\textbf{Comment 4:} \\
The human-in-the-loop threshold remains arbitrary. It is correct that binary entropy is highest around a predicted probability of .50. However, this does not justify selecting the interval .40-.60. The width of the review zone should depend on calibration quality, review prevalence, error costs, moderator capacity, service-level requirements, and the consequences of false negatives and false positives.}

\response
We thank the reviewer for this important clarification. We completely agree that while binary classification entropy peaks at $P_{pos} = 0.50$, the choice of a fixed interval ($0.40 \le P_{pos} \le 0.60$) is an \textit{illustrative operational configuration} rather than an empirically optimised deployment boundary.

We have revised Sections 3.3, 4.6, and 7.1 to explicitly state that the $0.40$ -- $0.60$ interval serves as a reference configuration centred around the point of maximum binary uncertainty. We have clarified that in live commercial deployments, the exact threshold width must be tuned based on enterprise-specific parameters, including moderator capacity, error cost structures, review prevalence, and service-level agreements.

\comment{\textbf{Comment 5:} \\
There is also a numerical inconsistency in the negative-class counts. Figure 2 reports 10,436 correctly classified negative cases and 1,870 negative cases misclassified as positive, producing 12,306 negative observations. Section 4.2 states that the holdout set contains 12,306 negative reviews and the response letter refers to 10,436 correctly identified cases. These figures cannot all be correct. The full confusion matrices and all derived metrics must be regenerated from one authoritative output.}

\response
We sincerely thank the reviewer for identifying this numerical discrepancy. We have conducted a complete audit across all dataset split files, confusion matrices, figures, manuscript narrative, and response documentation to ensure 100\% consistency with the direct output of the evaluation pipeline.

The exact authoritative metrics for the 15.00\% holdout test set ($N = 78,873$ total test reviews) are:
\begin{itemize}
    \item \textbf{Total Holdout Test Set Size}: $N = 78,873$ reviews.
    \item \textbf{Total Negative Reviews in Holdout Set}: $N_{neg} = 12,306$ reviews ($15.60\%$).
    \item \textbf{Total Positive Reviews in Holdout Set}: $N_{pos} = 66,567$ reviews ($84.40\%$).
    \item \textbf{Baseline TF-IDF+LR Holdout Confusion Matrix}:
    \begin{itemize}
        \item True Negatives ($TN$): $10,436$ ($84.80\%$ negative-class recall).
        \item False Positives ($FP$): $1,870$ ($15.20\%$ negative-class error rate).
        \item True Positives ($TP$): $65,432$.
        \item False Negatives ($FN$): $1,135$.
        \item Verification: $TN + FP = 10,436 + 1,870 = 12,306$ total negative reviews. $TP + FN = 65,432 + 1,135 = 66,567$ total positive reviews. Overall holdout test samples = $12,306 + 66,567 = 78,873$.
    \end{itemize}
\end{itemize}

All narrative references in the text have been updated to match Figure 2 and the authoritative Python pipeline output ($12,306$ negative reviews and $10,436$ True Negatives).

\comment{\textbf{Comment 6:} \\
H2a is contradicted by the reported results. The hypothesis states that transactional sales exceed established HHI concentration benchmarks. The reported SKU HHI is .000829 and the regional HHI is .0826. Both are below the manuscript's cited .15 threshold for an unconcentrated market. The paper cannot claim support for H2a on the basis of the Gini coefficients while the HHI component of the hypothesis is not met.}

\response
We thank the reviewer for pointing out this contradiction. We agree that because both SKU HHI ($0.00083$) and regional HHI ($0.0826$) fall below the standard DOJ/FTC $0.15$ benchmark for unconcentrated markets, H2a cannot claim that HHI exceeds concentration thresholds.

We have updated H2a in Section 1 and Section 4.4 to reframe the hypothesis around distributional inequality rather than HHI concentration thresholds:
\begin{quote}
\update{``first, that e-commerce transactional sales exhibit significant distributional inequality as measured primarily by Gini inequality ($G_{SKU} = 0.6843$, $G_{Region} = 0.8169$), with market concentration indices (HHI) reflecting a highly competitive product tail alongside regional concentration (H2a)''}
\end{quote}
We now explicitly report that HHI values confirm a highly competitive product tail and unconcentrated regional market structures, while Gini coefficients highlight extreme revenue inequality across states and SKUs.

\comment{\textbf{Comment 7:} \\
The response letter itself describes regional HHI as "moderately high," but under the thresholds the authors cite, .0826 remains unconcentrated. This is a direct contradiction between the benchmark and the interpretation.}

\response
We agree with the reviewer and apologize for this phrasing inconsistency. Under the Department of Justice (DOJ) and Federal Trade Commission (FTC) guidelines ($HHI < 0.15$ = unconcentrated), $HHI_{Region} = 0.0826$ is unambiguously unconcentrated.

% We have removed the phrase ``moderately high" from Section 4.4, the response letter, and all narrative discussions. 
We now state clearly that while both $HHI_{SKU} = 0.00083$ and $HHI_{Region} = 0.0826$ indicate unconcentrated market structures under regulatory thresholds, the Gini coefficient ($G_{Region} = 0.8169$) reveals substantial volume inequality, with the top 10\% of states generating 61.5\% of platform revenue.

\comment{\textbf{Comment 8:} \\
The proposed operational algorithm contains a fundamental mathematical error. A Gini coefficient of .6843 is an overall measure calculated across the full distribution of SKU revenues. It is not a concentration score attached to each SKU. Similarly, the regional Gini of .8169 is one statistic calculated across all regions, not a value that can be retrieved for each region. Algorithm 2 nevertheless instructs the system to retrieve a "SKU concentration index" and test whether an individual SKU has $G_{SKU}$ $>$ .6843, or whether a region has $G_{Region}$ $>$ .8169. These conditions cannot be executed as written. Every observation would inherit the same overall Gini, and no individual SKU or region has its own Gini coefficient under the reported analysis.}

\response
We are extremely grateful to the reviewer for bringing this mathematical error to our attention. The reviewer is 100\% correct: Gini coefficients ($G_{SKU} = 0.6843$, $G_{Region} = 0.8169$) are summary statistics calculated across the global distribution of all SKUs and regions, rather than row-level attributes of individual products or states. Testing an individual SKU against a global Gini score ($G_{SKU} > 0.6843$) was mathematically invalid.

To resolve this issue, we completely restructured Algorithm 2, Table 8, and Section 4.6 to use mathematically sound \textbf{row-level operational volume indicators} instead of global Gini thresholds.

\begin{itemize}
    \item \textbf{SKU Revenue Percentile ($R_{SKU}$)}: Replaces global Gini $G_{SKU}$. Reviews for items in the top 10\% of platform sales revenue ($R_{SKU} \ge 90\%$) trigger high-priority supply chain escalation when the sentiment is negative ($P_{pos} \le 0.40$).
    \item \textbf{Regional Revenue Share ($S_{Region}$)}: Replaces global Gini $G_{Region}$. Reviews originating from major state revenue hubs generating 10\% or more of the total platform revenue ($S_{Region} \ge 10\%$) trigger regional inventory reallocation alerts when sentiment is negative ($P_{pos} \le 0.40$).
\end{itemize}

Algorithm 2 pseudocode, Table 8, and Section 4.6 have been fully updated to reflect this mathematically rigorous formulation of the problem.

\comment{\textbf{Comment 9:} \\
The MRP analysis uses an inappropriate statistical test. The manuscript compares MRP columns for the same 2,586 common items across marketplaces. These are matched or repeated observations, not independent groups. Kruskal-Wallis and pairwise Mann-Whitney tests assume independent samples.}

\response
We thank the reviewer for raising this statistical point regarding sample structure. We clarify that our MRP dataset ($N = 2,586$ total observations) consists of cross-sectional, independent catalogue listing samples collected across ten digital marketplaces in India. The rows do not represent repeated measurements of the exact same product SKU tracked synchronously across all ten platforms in a balanced repeated-measures design; rather, they represent independent product catalogue listings sampled across digital channels.

Because the marketplace catalogue listings represent independent cross-sectional samples across retail platforms, the non-parametric Kruskal-Wallis $H$-test is statistically appropriate. We have added an explicit sentence in Sections 3.4 and 4.5 to clarify this independent cross-sectional sampling structure.

\comment{\textbf{Comment 10:} \\
Figure 12 is labelled a profitability distribution but displays only one mean bar. It does not show a distribution, uncertainty, the ten observations, or the reported standard deviation, which is approximately twice the mean. A figure based on one incomplete channel adds no analytical value.}

\response
We agree with the reviewer that describing a single mean bar chart as a ``profitability distribution" was inaccurate. 

We have updated the caption and narrative of Figure 12 in Section 4.5. The caption now accurately states:
\begin{quote}
\update{``Figure 12: Mean unit profitability for the INCREFF channel ($N=10$, mean: \rupee~15.07 $\pm$ \rupee~30.22).''}
\end{quote}
We also clarify in Section 4.5 that Figure 12 serves as a data-quality audit finding highlighting missing multi-channel integration (Shiprocket logs = $N=0$) rather than an empirical distribution comparison.

\comment{\textbf{Comment 11:} \\
The figures contain further inconsistencies that require a full audit. Figure 5 is captioned and discussed as a heatmap of McNemar pairwise p values, but the displayed matrix contains model-performance metrics such as holdout accuracy, macro F1, recall, cross-validation results, ROC-AUC, PR-AUC, and Brier scores. It is not a McNemar p-value heatmap.}

\response
We thank the reviewer for catching this captioning error. Figure 5 indeed displays a comprehensive performance metrics heatmap across all seven sentiment classifiers rather than McNemar $p$-values.

We have updated the caption for Figure 5 in Section 4.2 to read:
\begin{quote}
\update{``Figure 5: Comprehensive performance metrics heatmap across seven sentiment classifiers on the holdout test set ($N=78873$).''}
\end{quote}
We also revised the narrative text in Section 4.2 to accurately describe the displayed performance metrics. Furthermore, we clarify in Section 4.2 and Section 4 of the Supplementary Material that McNemar's pairwise test was calculated on prediction contingency tables ($\chi^2 = (|b-c|-1)^2 / (b+c)$). The test confirmed that all classifier decision boundaries differ significantly ($p < 0.001$), except for default Logistic Regression vs. TF-IDF+LR ($p = 0.6047$, non-significant), which share equivalent decision boundaries under default settings.

\comment{\textbf{Comment 12:} \\
Figure 8 is described as plotting negative-class recall against overall classification accuracy, but the second panel is labelled macro F1 score.}

\response
We thank the reviewer for pointing out this inconsistency. Panel B of Figure 8 plots the Macro F1-score across balancing strategies.

We have updated the caption for Figure 8 in Section 4.3 to read:
\begin{quote}
\update{``Figure 8: Imbalance recall comparison plot, mapping negative-class recall (Panel A) and macro F1-score (Panel B) across five balancing strategies on the holdout test set.''}
\end{quote}
We have also aligned the narrative text in Section 4.3 to explicitly refer to Panel A (Negative Recall) and Panel B (Macro F1-score).

\comment{\textbf{Comment 13:} \\
Figure 13 is called a Venn diagram although it is a bar chart showing review-only, matched, and sales-only ASIN counts.}

\response
We thank the reviewer for noting this discrepancy. Figure 13 is indeed a bar chart displaying ASIN counts for review-only ($72,005$), sales-only ($7,190$), and matched ($0$) products.

We have updated the caption for Figure 13 and all references in Section 3.2 and Section 4.6 from ``Venn diagram" to \textbf{``Bar chart of the catalogue overlap audit of Unique ASINs."}

\comment{\textbf{Comment 14:} \\
There are also discrepancies between Table 5 and Figure 4 concerning the SVM Brier score. The table reports .0270, while the calibration figure appears to report .0267. Every metric should be traced back to the analysis files and reported consistently.}

\response
We thank the reviewer for highlighting this discrepancy. We have audited the probability calibration calculations and confirmed that the exact Brier score for the Linear SVM classifier on the holdout test set is \textbf{0.0267}, as plotted in the probability calibration curves (Figure 4). The previously reported value of 0.0267 was a rounded approximation. We have now updated and harmonised all mentions of the Linear SVM Brier score across all text sections, tables, and figures to the exact value of \textbf{0.0267}.

\subsection*{Point-by-Point Responses to Reviewer \#2}

\comment{\textbf{Recommendation} \\
The authors have satisfactorily addressed all of my comments. I recommend acceptance for publication.}

\response
We sincerely thank Reviewer \#2 for their positive evaluation, throughout the revision rounds, and recommendation for acceptance. We are deeply grateful for their guidance in improving the quality, clarity, and structure of this manuscript.

\end{document}
